\subsubsection{Step Budget}\label{sec:res:fewstep}

The flow-based models keep their success with constraint satisfaction as the step budget falls to a
single network evaluation, while the diffusion baseline does not reach it at one evaluation
(\autoref{tab:k-ladder}). The baseline reaches $0.667$ at $\nfe=1$ and $1.000$ at $\nfe=2$, $10$ and
$20$. Over the same range the flow-based models never leave the top of the column: both average-velocity
models read $1.000$ or $0.967$ at one and two evaluations, analytic average-velocity matching holds
$0.933$ out to ten, and instantaneous-velocity matching reads $1.000$ at one, two and twenty.

\begin{table}[htpb]
  \caption[Success against step budget]{D3IL-avoiding, every model at every step budget of the grid
  $\nfe\in\{1,2,5,10,20\}$ under the per-step projector of \ac{DPCC}, each at the selection rule it is
  operated with: temporal consistency for the two average-velocity models, cumulative projection cost
  for instantaneous-velocity matching and for the baseline. The protocol of \ac{DPCC}: five training
  seeds $\times$ two episodes per geometry, averaged over the three geometries, tightened constraints,
  temporal U-Net of $4.0$\,M parameters; columns as in \autoref{tab:avoiding-dpcc-protocol}, mean $\pm$
  the standard deviation over the seeds. The other two rules of every cell are in
  \autoref{tab:app:avoiding-dpcc-full}. \emph{Pending}: cells not yet evaluated (R41); the diffusion
  baseline at $\nfe=5$ is a separately trained checkpoint. \selectedmark\ the operating point,
  \baselinemark\ the baseline of record.}
  \label{tab:k-ladder}
  \centering
  \footnotesize
  \begin{tabularx}{\linewidth}{@{}L c rrrrr@{}}
    \toprule
    Model (rule) & Column & $\nfe=1$ & $\nfe=2$ & $\nfe=5$ & $\nfe=10$ & $\nfe=20$ \\
    \midrule
    \selectedmark\ MeanFM ($t$) & S\&C & $0.967 \pm 0.075$ & $0.967 \pm 0.075$ & $0.933 \pm 0.149$ & $0.933 \pm 0.149$ & \emph{pending} \\
                & Steps   & $58.6 \pm 2.2$ & $59.4 \pm 2.0$ & $60.6 \pm 2.4$ & $63.6 \pm 1.9$ & \emph{pending} \\
                & ms/step & $18.3 \pm 0.2$ & $27.7 \pm 0.5$ & $223.8 \pm 14.5$ & $402.3 \pm 29.0$ & \emph{pending} \\
    \midrule
    CI-MeanFM, $\alpha_{\mathrm{end}}=0.2$ ($t$) & S\&C & $1.000 \pm 0.000$ & $1.000 \pm 0.000$ & \emph{pending} & \emph{pending} & \emph{pending} \\
                & Steps   & $59.2 \pm 1.2$ & $60.1 \pm 3.4$ & \emph{pending} & \emph{pending} & \emph{pending} \\
                & ms/step & $18.1 \pm 0.2$ & $27.0 \pm 0.2$ & \emph{pending} & \emph{pending} & \emph{pending} \\
    \midrule
    FM ($c$)    & S\&C & $1.000 \pm 0.000$ & $1.000 \pm 0.000$ & \emph{pending} & \emph{pending} & $1.000 \pm 0.000$ \\
                & Steps   & $67.0 \pm 0.8$ & $67.0 \pm 3.8$ & \emph{pending} & \emph{pending} & $63.2 \pm 1.5$ \\
                & ms/step & $17.3 \pm 0.1$ & $25.5 \pm 0.2$ & \emph{pending} & \emph{pending} & $476.7 \pm 26.0$ \\
    \midrule
    \baselinemark\ Diffusion ($c$) & S\&C & $0.667 \pm 0.391$ & $1.000 \pm 0.000$ & \emph{pending} & $1.000 \pm 0.000$ & $1.000 \pm 0.000$ \\
                & Steps   & $72.1 \pm 17.6$ & $61.2 \pm 1.2$ & \emph{pending} & $70.3 \pm 11.5$ & $70.1 \pm 8.2$ \\
                & ms/step & $33.3 \pm 14.5$ & $211.4 \pm 18.8$ & \emph{pending} & $309.8 \pm 11.2$ & $553.4 \pm 28.0$ \\
    \bottomrule
  \end{tabularx}
\end{table}
\dataref{Every printed cell is a row of \autoref{tab:app:avoiding-dpcc-full} (19-09 corpus; diffusion
$\nfe=2$ from \texttt{batch\_avoiding\_combined\_20260923\_110010}). The pending cells are R41 in
\texttt{data\_status/PENDING\_20260922\_all\_lacking\_runs.md}: MeanFM $\nfe=20$, CI-MeanFM $\nfe=5,10,20$
and FM $\nfe=5,10$ are evaluations of the existing checkpoints; diffusion $\nfe=5$ needs five trainings.
The earlier figure of this subsection (\texttt{fig\_avoiding\_k\_ladder}, drawn on ragged budgets) is
withdrawn: \texttt{withheld/20260923\_v3.69\_archive} (author, 23-09).}

The two halves of the table are not read the same way, and the difference is the point. Each diffusion
cell is its own model: the number of denoising steps is fixed when the noise schedule is discretised
during training, so the budgets of that row are separately trained checkpoints, one per step count, and
moving along it costs a training run. A flow-based row is one checkpoint per model, read at every
budget, because there the step count is the number of steps the ODE solver takes at inference. For a
flow-based model the budget is therefore a dial that can be turned after training, and for the diffusion
baseline it is a property of the checkpoint --- the same asymmetry \autoref{fig:raw-plans} shows, where
the two diffusion panels are two separately trained models and each flow-based row is one model read
twice. The cost of that asymmetry is a training run per budget, and what it buys is visible in the last
row: at two denoising steps the baseline holds $1.000$ at $211.4$\,ms per control step, against
instantaneous-velocity matching's $17.3$\,ms at one evaluation. Cutting the baseline's budget from twenty
to two leaves it at twelve times the cost, most of it in the projection of its first denoising step
(\autoref{sec:res:avoiding:caveat:noise}).

